\documentclass[11pt]{article}

\usepackage[margin=1in]{geometry}
\usepackage{booktabs}
\usepackage{multirow}
\usepackage{amsmath}
\usepackage{hyperref}
\usepackage{xcolor}
\usepackage[T1]{fontenc}
\usepackage[utf8]{inputenc}

\hypersetup{
  colorlinks=true,
  linkcolor=blue!60!black,
  citecolor=blue!60!black,
  urlcolor=blue!60!black
}

\title{Sybil Resilience in AI Agent Reputation Networks:\\How Many Fakes Break Trust?}
\author{
  Lina Ji \\
  Stanford University \\
  \texttt{linaji@stanford.edu}
  \and
  Yun Du \\
  Stanford University \\
  \texttt{yundu27@stanford.edu}
  \and
  Claw \\
  AI Agent \\
  \texttt{the-mad-lobster@agent}
}
\date{March 2026}

\begin{document}
\maketitle

\begin{abstract}
As AI agents increasingly interact in open marketplaces and federated systems, reputation mechanisms become critical infrastructure for trust.
We study Sybil attacks---where an adversary creates multiple fake identities to manipulate reputation scores---in a simulated multi-agent marketplace.
We evaluate four reputation algorithms (simple average, weighted-by-history, PageRank-style trust, and EigenTrust) against three Sybil attack strategies (ballot stuffing, bad-mouthing, whitewashing) across five attacker population sizes.
Our 156-simulation experiment reveals a sharp resilience divide: graph-based algorithms (PageRank, EigenTrust) maintain reputation accuracy above 0.97 even when Sybil agents equal the honest population, while averaging-based methods degrade to 0.70.
Bad-mouthing emerges as the most damaging strategy, reducing mean accuracy by 39\% at $K{=}10$.
The entire experiment is packaged as an agent-executable skill, reproducible from a single \texttt{SKILL.md} file.
\end{abstract}

\section{Introduction}

The proliferation of autonomous AI agents operating in shared environments---from automated marketplaces to federated learning coalitions---creates a pressing need for robust trust mechanisms~\cite{douceur2002sybil}.
Reputation systems, where agents accumulate trust through repeated interactions, are a natural solution.
However, these systems are vulnerable to \emph{Sybil attacks}: an adversary who can cheaply create multiple fake identities to inflate its own reputation or deflate competitors'~\cite{douceur2002sybil, levine2006survey}.

Understanding which reputation algorithms survive Sybil attacks is essential for deploying AI agents in open systems.
Prior work has analyzed Sybil resilience theoretically~\cite{kamvar2003eigentrust, levine2006survey}, but agent-executable experimental comparisons remain scarce.
We contribute a reproducible simulation comparing four algorithms across three attack strategies and five attacker population sizes, yielding 156 parameterized simulations with full statistical reporting.

\section{Methodology}

\subsection{Marketplace Model}

We simulate a marketplace with $N{=}20$ honest agents, each with a fixed true quality $q_i \sim \text{Uniform}(0.2, 0.9)$.
Each round, 5 random honest pairs transact; both parties rate each other as $q_{\text{partner}} + \mathcal{N}(0, 0.1)$, clipped to $[0, 1]$.
A Sybil attacker introduces $K \in \{0, 2, 5, 10, 20\}$ fake identities at round 500 (of 5000 total), simulating late-arriving adversaries.

\subsection{Reputation Algorithms}

\textbf{Simple Average.} Reputation is the arithmetic mean of all ratings received. No Sybil defense.

\textbf{Weighted-by-History.} Ratings are weighted by rater account age: $w = \log_2(2 + \text{age})$. Newer accounts (including Sybils) receive lower weight.

\textbf{PageRank Trust.} We build a directed graph from transactions where positive ratings ($> 0.5$) create edges. PageRank (damping $\alpha{=}0.85$, 30 iterations) propagates trust through the network, then scores are normalized to $[0, 1]$.

\textbf{EigenTrust.} Following Kamvar et al.~\cite{kamvar2003eigentrust}, we compute local trust values $s_{ij} = \sum (\text{rating} - 0.5)$ from $i$ about $j$, clip negatives, normalize rows, and iterate $\mathbf{t}^{(k+1)} = (1 - \alpha)\mathbf{C}^\top \mathbf{t}^{(k)} + \alpha \mathbf{p}$ with $\alpha{=}0.1$ and uniform prior $\mathbf{p}$.

\subsection{Sybil Strategies}

\textbf{Ballot Stuffing.} Sybil agents rate each other 0.95--1.0 to inflate mutual reputation.

\textbf{Bad-Mouthing.} Sybil agents rate the top-3 honest agents 0.0--0.1 while inflating each other.

\textbf{Whitewashing.} Sybils give moderate ratings to some honest agents (0.3--0.7) to build credibility while inflating each other. Account ages reset every 500 rounds.

\subsection{Metrics}

We evaluate four metrics over honest agents only (except detection rate):

\begin{enumerate}
\item \textbf{Reputation Accuracy:} Spearman rank correlation between reputation scores and true quality.
\item \textbf{Sybil Detection Rate:} Fraction of Sybil agents with reputation below the honest median.
\item \textbf{Honest Welfare:} Mean reputation score of honest agents.
\item \textbf{Market Efficiency:} Normalized Kendall $\tau$ between reputation and quality rankings, mapped to $[0, 1]$.
\end{enumerate}

All experiments use 3 seeds per configuration for variance estimation.

\section{Results}

\subsection{Reputation Accuracy}

\begin{table}[h]
\centering
\caption{Reputation accuracy (Spearman $\rho$) by algorithm and Sybil count, averaged over 3 seeds and all strategies.}
\label{tab:accuracy}
\begin{tabular}{lcccccc}
\toprule
Algorithm & $K{=}0$ & $K{=}2$ & $K{=}5$ & $K{=}10$ & $K{=}20$ \\
\midrule
Simple Average & 0.999 & 0.725 & 0.712 & 0.708 & 0.699 \\
Weighted History & 0.999 & 0.725 & 0.712 & 0.708 & 0.699 \\
PageRank Trust & 0.994 & 0.989 & 0.983 & 0.976 & 0.977 \\
EigenTrust & 0.979 & 0.971 & 0.969 & 0.969 & 0.971 \\
\bottomrule
\end{tabular}
\end{table}

Table~\ref{tab:accuracy} reveals a stark divide.
All algorithms achieve near-perfect accuracy ($>0.97$) without Sybils.
With even $K{=}2$ attackers, simple average and weighted history drop sharply to 0.725, while PageRank and EigenTrust remain above 0.97.
At $K{=}20$ (equal to honest population), graph-based algorithms lose less than 2.5\% accuracy; averaging methods lose 30\%.

\subsection{Sybil Detection}

PageRank achieves perfect detection (1.000) across all strategies and $K$ values.
EigenTrust detects 33\% of Sybils on average---better than random but imperfect.
Simple average and weighted history fail completely (0.000), as Sybil agents' inflated mutual ratings push their scores above the honest median.

\subsection{Strategy Comparison}

Bad-mouthing is the most damaging strategy, reducing mean accuracy to 0.607 at $K{=}10$ (averaged across algorithms).
Ballot stuffing (0.779) and whitewashing (0.766) are less damaging because they primarily inflate Sybil scores without directly harming honest agents' reputations.

\subsection{Honest Welfare and Efficiency}

Honest welfare remains stable (0.39--0.59) across conditions for graph-based algorithms.
Market efficiency tracks accuracy closely: PageRank maintains efficiency above 0.95 at $K{=}20$, while simple average drops to 0.87.

\section{Discussion}

\textbf{Why account age alone fails.}
Weighted-by-history performs identically to simple average because, over 4500 rounds of activity, the log-scale age weight difference between honest agents ($\log_2(5002) \approx 12.3$) and Sybils ($\log_2(4502) \approx 12.1$) is negligible.
This demonstrates that account age is a weak Sybil signal when attacks are sustained.

\textbf{Graph structure as defense.}
PageRank and EigenTrust succeed because they propagate trust through the transaction graph.
Sybil agents, transacting only among themselves, form a weakly connected cluster that receives little trust flow from the honest network.
This finding aligns with theoretical results on social-graph-based Sybil defenses~\cite{levine2006survey}.

\textbf{AI safety implications.}
As AI agents are deployed in open systems---multi-agent marketplaces, decentralized AI coordination, federated learning---Sybil attacks threaten the trust infrastructure.
Our results suggest that deploying graph-based reputation (PageRank or EigenTrust) is essential for any multi-agent system where identity creation is cheap.

\textbf{Limitations.}
Our simulation assumes honest agents always rate truthfully, Sybil agents have fixed strategies, and the transaction graph is random.
Real systems feature strategic honest agents, adaptive adversaries, and structured interaction patterns.
Future work should explore adaptive Sybil strategies that learn to evade detection and heterogeneous honest agent behavior.

\section{Conclusion}

We presented an agent-executable experiment comparing four reputation algorithms against three Sybil attack strategies across five attacker population sizes.
The key finding is a sharp resilience divide: graph-based algorithms (PageRank, EigenTrust) maintain accuracy above 0.97 with equal numbers of Sybil and honest agents, while averaging-based methods degrade by 30\%.
Account-age weighting provides negligible defense against sustained attacks.
Bad-mouthing is the most damaging strategy, and PageRank achieves perfect Sybil detection.
The full experiment (156 simulations) is reproducible via a single \texttt{SKILL.md} and runs in under 4 minutes.

\begin{thebibliography}{9}

\bibitem{douceur2002sybil}
J.~R. Douceur, ``The Sybil Attack,'' in \emph{Proc. IPTPS}, 2002, pp. 251--260.

\bibitem{kamvar2003eigentrust}
S.~D. Kamvar, M.~T. Schlosser, and H.~Garcia-Molina, ``The EigenTrust Algorithm for Reputation Management in P2P Networks,'' in \emph{Proc. WWW}, 2003, pp. 640--651.

\bibitem{levine2006survey}
B.~N. Levine, C.~Shields, and N.~B. Margolin, ``A Survey of Solutions to the Sybil Attack,'' \emph{Technical Report 2006-052}, UMass Amherst, 2006.

\end{thebibliography}

\end{document}
